\subsubsection*{Численный пример}

\begin{idea}[title={Откуда этот материал}]
На лекции метод Зойтендейка разбирался только теоретически. Ниже две итерации
на числах: видно, как вспомогательная задача выбирает направление вдоль
границы, пока ограничение активно, и переходит к антиградиенту, когда точка
уходит внутрь.
\end{idea}

\begin{example}\label{ex:zoutendijk-run}
\[
  \begin{cases}
    \varphi_0(x) = (x_1 - 2)^2 + (x_2 - 1)^2 \to \min , \\
    x_1 \ge 0 , \quad x_2 \ge 0 ,
  \end{cases}
\]
то есть в стандартной форме
\[
  \varphi_1(x) = -x_1 \le 0 , \qquad \varphi_2(x) = -x_2 \le 0 .
\]
Ограничений-равенств нет, поэтому условие $As = 0$ отсутствует. Берём
$\xi_0 = \xi_1 = \xi_2 = 1$, $\delta_0 = 0{,}1$. Градиенты:
\[
  \nabla\varphi_0(x) = \begin{pmatrix} 2(x_1-2) \\ 2(x_2-1) \end{pmatrix} ,
  \qquad
  \nabla\varphi_1 = \begin{pmatrix} -1 \\ 0 \end{pmatrix} ,
  \qquad
  \nabla\varphi_2 = \begin{pmatrix} 0 \\ -1 \end{pmatrix} .
\]
Старт: $x_0 = (0, 2)\T \in S$, $\varphi_0(x_0) = 4 + 1 = 5$.

\medskip
\noindent\textbf{Итерация 0.} Значения ограничений:
$\varphi_1(x_0) = 0$, $\varphi_2(x_0) = -2$. Первое активно, второе даже не
почти активно ($-2 < -\delta_0$), значит
$I_{\delta_0}(x_0) = \{1\}$. Градиент цели
$\nabla\varphi_0(x_0) = (-4, 2)\T$.

Вспомогательная задача:
\[
  \eta \to \min
  \qquad\text{при}\qquad
  \begin{cases}
    -4 s^{(1)} + 2 s^{(2)} \le \eta , \\
    -s^{(1)} \le \eta , \\
    \abs{s^{(1)}} \le 1 , \quad \abs{s^{(2)}} \le 1 .
  \end{cases}
\]
Из второго ограничения $\eta \ge -s^{(1)} \ge -1$, то есть $\eta$ не может быть
меньше $-1$. Значение $\eta_0 = -1$ достижимо: при $s^{(1)} = 1$ второе условие
обращается в равенство, а первое принимает вид $2s^{(2)} \le 3$ и выполнено при
любом допустимом $s^{(2)}$. Оптимум вырожден — годится целое ребро; берём
вершину
\[
  s_0 = \begin{pmatrix} 1 \\ -1 \end{pmatrix} ,
  \qquad \eta_0 = -1 < 0 ,
\]
дающую наибольшее убывание цели:
$\nabla\T\varphi_0(x_0)\,s_0 = -4 - 2 = -6 < 0$.

Шаг. Точка $x(\alpha) = (\alpha,\ 2-\alpha)\T$ допустима при
$0 \le \alpha \le 2$, и
\[
  \varphi_0\big(x(\alpha)\big) = (\alpha - 2)^2 + (1 - \alpha)^2
    = 2\alpha^2 - 6\alpha + 5 ,
  \qquad
  \varphi_0' = 4\alpha - 6 = 0 \implies \alpha_0 = \tfrac32 \in [0,2] .
\]
\[
  x_1 = \begin{pmatrix} 3/2 \\ 1/2 \end{pmatrix} ,
  \qquad
  \varphi_0(x_1) = \tfrac14 + \tfrac14 = \tfrac12
  \qquad \big(5 \to \tfrac12\big) .
\]

\medskip
\noindent\textbf{Итерация 1.} Теперь $\varphi_1(x_1) = -\tfrac32$,
$\varphi_2(x_1) = -\tfrac12$ — обе строго меньше $-\delta_0$, значит
$I_{\delta_0}(x_1) = \varnothing$: точка ушла внутрь допустимого множества.
Градиент $\nabla\varphi_0(x_1) = (-1, -1)\T$, вспомогательная задача
вырождается в
\[
  \eta \to \min
  \qquad\text{при}\qquad
  -s^{(1)} - s^{(2)} \le \eta , \quad \abs{s^{(i)}} \le 1 ,
\]
откуда $s_1 = (1, 1)\T$, $\eta_1 = -2$. Это ровно антиградиент, нормированный
условием $\norm{s}_\infty \le 1$: ограничения на направление не влияют, и метод
работает как в безусловной задаче.

Шаг. $x(\alpha) = \big(\tfrac32 + \alpha,\ \tfrac12 + \alpha\big)\T$ допустима
при всех $\alpha \ge 0$, и
\[
  \varphi_0\big(x(\alpha)\big)
   = \Big(\alpha - \tfrac12\Big)^2 + \Big(\alpha - \tfrac12\Big)^2
   = 2\Big(\alpha - \tfrac12\Big)^2
  \quad\Longrightarrow\quad \alpha_1 = \tfrac12 ,
\]
\[
  x_2 = \begin{pmatrix} 2 \\ 1 \end{pmatrix} = x_* ,
  \qquad \varphi_0(x_2) = 0
  \qquad \big(\tfrac12 \to 0\big) .
\]
Значения цели: $5 \to \tfrac12 \to 0$.
\end{example}

\begin{idea}
Что здесь существенно.

\begin{itemize}[nosep,left=0pt]
  \item На нулевой итерации чистый антиградиент
        $-\nabla\varphi_0(x_0) = (4,-2)\T$ вывел бы точку за границу
        $x_2 \ge 0$ только при большом шаге, но главное другое: в вершине
        $x_1 = 0$ он смотрит \emph{наружу} по первой координате. Условие
        $\nabla\T\varphi_1 s \le \xi_1 \eta < 0$ заставляет брать
        $s^{(1)} > 0$ — направление отталкивается от активного ограничения.
  \item Вырожденность оптимума вспомогательной задачи — типичная ситуация:
        $\eta$ упирается в самое «тугое» ограничение, а остальные координаты
        направления остаются недоопределёнными. Симплекс-метод вернёт
        какую-то вершину; все они дают $\eta_0 = -1$, но убывание цели вдоль
        них разное.
  \item Шаг здесь брался точной минимизацией, а не по оценкам
        \eqref{eq:zout-6}--\eqref{eq:zout-8} из леммы~\ref{lem:zoutendijk-step}:
        те оценки нужны для доказательства сходимости, на практике же
        достаточно убедиться, что точка осталась допустимой и цель убыла.
\end{itemize}
\end{idea}
